\documentclass[UTF8,zihao=-4]{ctexart}
\usepackage[a4paper,top=2.5cm,bottom=2.5cm,left=2.8cm,right=2.8cm]{geometry}
\usepackage{array}
\usepackage{booktabs}
\usepackage{enumitem}

\setlength{\parindent}{2em}
\setlength{\parskip}{0.35em}
\linespread{1.25}
\pagestyle{plain}

\begin{document}

\begin{center}
    {\zihao{2}\bfseries 本科毕业设计（论文）修改情况说明}
\end{center}

\vspace{1em}

\begin{tabular}{>{\bfseries}p{3cm}p{11cm}}
论文题目： & 基于动作扩散策略的四足机器人视觉目标导航 \\
学生姓名： & 陈样 \\
学号： & 522031910056 \\
指导教师： & 雷华明 \\
\end{tabular}

\section*{一、修改说明}

根据评阅人和答辩小组提出的意见，本人已对毕业设计（论文）进行了认真检查和修改。修改内容主要包括：补充研究动机与国内外研究现状，精简部分基础理论介绍，进一步明确论文研究对象与平台部署定位，补充真机实验到达判定条件，完善 TTS 实验综合评分函数及参数说明，规范摘要关键词表述，并同步更新英文大摘要与电子归档文件。现将主要修改情况说明如下。

\section*{二、评阅与答辩意见及修改情况}

\subsection*{1. 关于动作扩散机制引入动机不够明确的问题}

\noindent\textbf{修改建议：}论文对动作扩散机制在四足机器人上部署研究不足的原因分析不够深入，因此引入动作扩散机制的动机不够明确，建议进行简要补充。

\noindent\textbf{修改情况：}已在论文绪论及相关研究部分补充动作扩散策略、视觉导航和四足机器人真实部署相关研究现状，进一步说明现有工作在边缘端实时推理、四足平台闭环部署和真实系统验证方面仍存在不足。相应内容强化了本文选择动作扩散策略作为高层视觉导航方法，并将其部署到四足机器人平台上的研究动机。

\subsection*{2. 关于第二章内容偏多、基础介绍需精简的问题}

\noindent\textbf{修改建议：}第二章内容较多，建议精简前面几节关于通用模型介绍的部分，突出作者自己的工作和内容。

\noindent\textbf{修改情况：}已对第二章中扩散模型、DDPM 和 DDIM 采样机制等基础理论介绍进行压缩和整合，减少与本文实验关系不够直接的通用性叙述。同时保留与本文统一推理模块、DDIM 加速、CFG 引导、TTS 候选筛选和视觉编码器替换直接相关的理论与实验内容，使章节重点更加集中于本文完成的算法设计、实验分析和部署配置选择。

\subsection*{3. 关于论文研究对象表述的问题}

\noindent\textbf{修改建议：}论文不应仅表述为四足机器人相关工作，也应突出其作为一种视觉导航方法的研究属性。

\noindent\textbf{修改情况：}已对论文中的研究对象表述进行调整，明确本文研究的是面向四足机器人平台部署的视觉目标导航方法。论文既强调算法层面的动作扩散视觉导航与候选轨迹生成，也说明四足机器人平台在真实部署中的相机扰动、速度接口、安全状态机和闭环频率等特殊约束，从而兼顾“视觉导航方法”和“四足机器人真实部署”两个层面。

\subsection*{4. 关于真机实验到达判定条件的问题}

\noindent\textbf{修改建议：}仿真实验中设置机器人距离目标 0.5 m 以内视为到达，但现实世界实验中没有通过外部测量方式严格判断到达，建议补充说明。

\noindent\textbf{修改情况：}已在真机实验章节补充真实平台上的到达判定条件。论文中明确说明，MuJoCo 仿真能够读取机器人位姿和目标点坐标，因此采用 0.5 m 的物理距离阈值作为成功判据；而真机实验未引入外部定位或运动捕捉系统，因此采用视觉拓扑到达判据，即当实时子目标节点达到或超过任务目标节点，且 NoMaD 距离预测头输出的视觉拓扑距离小于阈值时，认为机器人到达目标区域。论文同时说明该阈值不是米制距离，而是模型距离预测输出空间中的无量纲阈值。

\section*{三、其他修改情况}

除上述针对评阅与答辩意见所作的修改外，论文还进行了以下完善：

\begin{enumerate}[label=\arabic*.]
    \item 补充 TTS 实验中候选轨迹评分函数的具体表达式及权重参数，明确前向进展、横向偏摆和平滑性代价在候选筛选中的作用。
    \item 在 TTS 单因素实验和 DDIM、CFG、TTS 联合实验表格附近补充配置级综合评分函数，说明综合评分由 Forward Progress、Latency、Lateral Abs 和 Smoothness 加权计算得到，并给出具体权重。
    \item 对论文部分章节的表述进行语言润色，使研究动机、实验设置、指标含义和结论分析更加清晰。
    \item 更新英文大摘要，使其与最新表述保持一致。
    \item 规范中文摘要关键词，将中文关键词压缩为 5 个中文技术词条，并在英文关键词中保留 NoMaD、DDIM 等通用英文缩写以便检索。
\end{enumerate}

\section*{四、导师确认}

本人已根据评阅人和答辩小组意见完成论文修改，并提交指导教师审阅确认。

\vspace{2em}

\begin{tabular}{p{6cm}p{6cm}}
学生签名： & 指导教师签名： \\
\vspace{2em} & \vspace{2em} \\
日期：2026 年5月26 日 & 日期：2026 年5 月26 日 \\
\end{tabular}

\end{document}
